% !PRINTED 14
\weberpage{14}
%<<CONT
manifold and have the same arbitrarily chosen denominator. Addition is
then defined so that
\[
 {a\over c}+{b\over c}={a+b\over c}
\]
holds.
This rule contains as a special case the addition of rational
fractions. To justify it generally it remains only to show that, if
$a:c=a':c'$ and $b:c=b':c'$, then also
$(a+b):c=(a'+b'):c'$ must hold. We prove this by showing that, if
$a:c=a':c'$ and $(a+b):c>(a'+b'):c'$, then also $b:c>b':c'$ must
hold. Let therefore, with $m,n$ two whole numbers,
\[
 {a+b\over c}>{m\over n}>{a'+b'\over c'},
\]
then $n(a+b)>mc$, and hence
\[
 {b\over c}>{mc-na\over nc},
 \qquad {mc'-na'\over nc'}>{b'\over c'}.
\]
On the other hand, from the supposition $a:c=a':c'$ it follows easily
that
\[
 {mc-na\over nc}={mc'-na'\over nc'}
\]
holds, and therefore $b:c>b':c'$, as was required to be proved.

It is thus shown that the numbers, as we have defined them, also form a
measurable set. If $a$ and $c$ are taken from a continuous manifold,
then, with fixed $c$, the ratios $a:c$ also form a continuous set;
hence, since continuous sets exist at all, it follows that the numbers
too form a continuous set.

If $\alpha,\beta,\gamma,\delta$ are now numbers, then from the
proportion $\alpha:\beta=\gamma:\delta$ any one of the four numbers
can be determined by the three others given. If one puts $\delta=1$
and seeks $\alpha$, one obtains the \gesperrt{multiplication}
$\alpha=\beta\gamma$; the commutativity of the factors is a consequence
of the theorem that from $\alpha:\beta=\gamma:\delta$ there follows
$\alpha:\gamma=\beta:\delta$. If one seeks $\gamma$, one obtains
division; and from the definition of addition given above
there follows the fundamental formula
$\alpha(\beta+\gamma)=\alpha\beta+\alpha\gamma$. The four fundamental
operations are therefore executable in the domain of numbers, with the
sole restriction that in subtraction the subtrahend must be smaller
than the minuend.

\begin{center}
\rule{0.18\linewidth}{0.4pt}
\end{center}
